% !TEX root = ../algebra_02.tex

\section{Условия, эквивалентные обратимости матрицы}


\begin{Theorem}{Критерии обратимости квадратной матрицы}{}
	Пусть $A \in M_n(K)$ --- квадратная матрица порядка $n$. Следующие условия эквивалентны:
	\begin{enumerate}
		\item Матрица $A$ обратима (существует $A^{-1} \in M_n(K)$ такая, что $A A^{-1} = A^{-1} A = I_n$).
		\item Ранг матрицы максимален: $\mathrm{rk}(A) = n$.
		\item Столбцы матрицы $A$ линейно независимы (образуют базис в $K^n$).
		\item Строки матрицы $A$ линейно независимы.
		\item Однородная система линейных уравнений $Ax = 0$ имеет единственное (тривиальное) решение $x = 0$.
		\item Линейное отображение $\mathcal{A} \colon K^n \to K^n$, заданное как $\mathcal{A}(x) = Ax$, является биективным (изоморфизмом).
		\item Определитель матрицы отличен от нуля: $\det A \neq 0$.
		\item Нуль не является собственным значением матрицы $A$.
	\end{enumerate}
\end{Theorem}

\begin{proof}
	(Некоторые утверждения в доказательстве будут доказаны дальше)
	\begin{itemize}
		\item \textbf{(2) $\Leftrightarrow$ (3) $\Leftrightarrow$ (4):} По определению столбцовый ранг --- это размерность линейной оболочки столбцов. Если ранг равен $n$, то $n$ столбцов образуют базис и линейно независимы. То же верно для строк в силу теоремы $\mathrm{rk}_{\mathrm{col}}(A) = \mathrm{rk}_{\mathrm{row}}(A)$.
		
		\item \textbf{(2) $\Leftrightarrow$ (5):} По теореме о множестве решений ОСЛУ $Ax = 0$, пространство решений является ядром $\Ker\mathcal{A}$ линейного оператора и имеет размерность $n - \mathrm{rk}(A)$. Единственность решения эквивалентна $\dim\Ker\mathcal{A} = 0$, что выполняется тогда и только тогда, когда $\mathrm{rk}(A) = n$.
		
		\item \textbf{(5) $\Leftrightarrow$ (6):} Тривиальность ядра $\Ker\mathcal{A} = \{0\}$ эквивалентна инъективности $\mathcal{A}$. Поскольку отображение действует из пространства в себя ($\dim = n < \infty$), по аналогу принципа Дирихле для линейных операторов инъективность равносильна сюръективности, а значит и биективности (изоморфизму).
		
		\item \textbf{(6) $\Rightarrow$ (1):} Если $\mathcal{A}$ --- изоморфизм, то существует обратное линейное отображение $\mathcal{A}^{-1}$. По теореме о матрице композиции, матрица отображения $\mathcal{A}^{-1}$ будет являться обратной матрицей к $A$.
		
		\item \textbf{(7) $\Leftrightarrow$ (8):} По определению характеристического многочлена $\chi_A(x) = \det(A - xI_n)$. Значение многочлена в нуле равно $\chi_A(0) = \det A$. Корни $\chi_A(x)$ --- это в точности собственные значения матрицы. Следовательно, $\lambda = 0$ не является собственным значением тогда и только тогда, когда $\det A \neq 0$.
	\end{itemize}
\end{proof}
